\section{Experiments}
\label{sec:experiments}

\subsection{Genre Classification}

\paragraph{Cross-validation.}
We evaluate both classifiers using stratified 5-fold cross-validation on the full corpus before committing to a final model. Each fold preserves the original class distribution. Table~\ref{tab:cv} summarizes the per-fold accuracy for both the genre and sentiment classifiers.

\begin{table}[h]
\centering
\begin{tabular}{lcc}
\toprule
\textbf{Fold} & \textbf{Genre Acc.} & \textbf{Sentiment Acc.} \\
\midrule
Fold 1 & 0.9685 & 0.7760 \\
Fold 2 & 0.9820 & 0.8009 \\
Fold 3 & 0.9730 & 0.7964 \\
Fold 4 & 0.9730 & 0.8118 \\
Fold 5 & 0.9775 & 0.7959 \\
\midrule
\textbf{Mean} & \textbf{0.9748} & \textbf{0.7962} \\
\bottomrule
\end{tabular}
\caption{Stratified 5-fold cross-validation accuracy for both classifiers. Sentiment CV uses BBC articles pseudo-labeled by VADER.}
\label{tab:cv}
\end{table}

\paragraph{Test set results.}
The final model is trained on the full training partition (80\%, $\approx$1,777 articles) and evaluated on the held-out test set (10\%, 227 articles). It achieves \textbf{98.24\% accuracy}. Per-class precision, recall, and F1 are shown in Table~\ref{tab:classification_report}.

\begin{table}[h]
\centering
\begin{tabular}{lcccc}
\toprule
\textbf{Category} & \textbf{Prec.} & \textbf{Recall} & \textbf{F1} & \textbf{Support} \\
\midrule
Business      & 1.0000 & 0.9608 & 0.9800 & 51 \\
Entertainment & 1.0000 & 0.9750 & 0.9873 & 40 \\
Politics      & 0.9556 & 1.0000 & 0.9773 & 43 \\
Sport         & 0.9811 & 1.0000 & 0.9905 & 52 \\
Tech          & 0.9756 & 0.9756 & 0.9756 & 41 \\
\midrule
\textbf{Weighted avg} & \textbf{0.9829} & \textbf{0.9824} & \textbf{0.9824} & \textbf{227} \\
\bottomrule
\end{tabular}
\caption{Per-class classification report on the 227-article test set.}
\label{tab:classification_report}
\end{table}

\paragraph{Confusion matrix.}
The confusion matrix (Table~\ref{tab:confusion}) reveals that the vast majority of errors are between adjacent topic areas — a business article misclassified as politics, or a sport article misclassified as tech. There are no gross misclassifications (e.g., sport predicted as entertainment).

\begin{table}[h]
\centering
\setlength{\tabcolsep}{4pt}
\begin{tabular}{lccccc}
\toprule
\textbf{True \textbackslash{} Pred} & \textbf{Bus.} & \textbf{Ent.} & \textbf{Pol.} & \textbf{Spt.} & \textbf{Tch.} \\
\midrule
Business      & \textbf{49} & 0 & 1 & 0 & 1 \\
Entertainment & 0 & \textbf{39} & 1 & 0 & 0 \\
Politics      & 0 & 0 & \textbf{43} & 0 & 0 \\
Sport         & 0 & 0 & 0 & \textbf{52} & 0 \\
Tech          & 0 & 0 & 0 & 1 & \textbf{40} \\
\bottomrule
\end{tabular}
\caption{Confusion matrix on the test set. Rows = true labels, columns = predicted. Diagonal entries (bold) are correct predictions.}
\label{tab:confusion}
\end{table}

% -----------------------------------------------------------------------
% TODO: A heatmap version of this confusion matrix (e.g., generated
%   with matplotlib/seaborn and exported as a PDF or PNG) would be
%   much easier to read than a plain table. Include it as a figure.
% -----------------------------------------------------------------------

\subsection{Keyword and Entity Extraction}

We present three sample articles from the corpus to illustrate the keyword and NER pipeline. Each article's top-8 TF-IDF keywords, top-8 POS-filtered keywords, and named entities are shown in Table~\ref{tab:samples}.

\begin{table*}[t]
\centering
\begin{tabular}{p{2.4cm}p{4.8cm}p{4.8cm}p{3.2cm}}
\toprule
\textbf{Article} & \textbf{TF-IDF Keywords} & \textbf{POS Keywords} & \textbf{Named Entities} \\
\midrule
\textit{Time Warner profits} (business)
  & timewarner, yearearlier, aols, aol, restate, fullyear, 639m, 111bn
  & timewarner, aol, profits, sales, warner, internet, profit, time
  & GPE: Ad, Google; PERSON: Time Warner, Warner; ORG: TimeWarner \\
\addlinespace
\textit{Dollar/Greenspan} (business)
  & greenspan, 12871, 12974, greenspans, sinche, loosening, halfpoint, deficit
  & dollar, deficit, months, greenspan, federal, reserve, current
  & PERSON: Alan Greenspan; ORG: Federal Reserve; FACILITY: White House \\
\addlinespace
\textit{Yukos/Rosneft} (business)
  & yugansk, rosneft, menatep, yukos, 479m, 540m, rosnefts, moscowbased
  & rosneft, yukos, sale, unit, loan, menatep, assets
  & GPE: Russian; ORG: Yukos, Menatep Group \\
\bottomrule
\end{tabular}
\caption{Sample keyword and entity extraction output for three business articles from the corpus.}
\label{tab:samples}
\end{table*}

Several observations stand out. TF-IDF keywords tend to capture proper nouns and rare compound forms (e.g., \textit{yearearlier}, \textit{moscowbased}) that are unique to the article, while POS-filtered keywords tend to be more general and readable. The NER output correctly identifies organizations and people but occasionally misclassifies common nouns as geopolitical entities (e.g., ``Ad'' as a GPE), a known limitation of NLTK's rule-based chunker.

\subsection{Sentiment Analysis}

Table~\ref{tab:sentiment} shows the sentiment predictions for the same three articles from both the NB model and VADER.

\begin{table}[h]
\centering
\begin{tabular}{lccc}
\toprule
\textbf{Article} & \textbf{NB} & \textbf{VADER} & \textbf{Compound} \\
\midrule
Time Warner profits & pos & pos & $+$0.9954 \\
Dollar/Greenspan    & neg & pos & $+$0.9421 \\
Yukos/Rosneft       & neg & neg & $-$0.8860 \\
\bottomrule
\end{tabular}
\caption{NB and VADER sentiment predictions with VADER compound scores for three sample articles. The second article shows a NB/VADER disagreement.}
\label{tab:sentiment}
\end{table}

Agreement between the two systems is mixed. On the three samples, two agree and one disagrees — a pattern we observe broadly across the corpus. We examine the causes of disagreement in Section~\ref{sec:discussion}.
